% year: תשפ"ד
% semester: א
% moed: א
% number: 5
% points: 14
% categories: antiderivatives
% source: exams/מבחנים/תשפד/מועד א תשפד.pdf

\begin{question}
חשבו את הפונקציה הקדומה של
\[ f(x)=\frac{4x^2+9x+10}{x^3+2x^2+5x} \]
\end{question}

\begin{hint}
פרקו את המכנה: $x^3+2x^2+5x=x(x^2+2x+5)$, כאשר לגורם הריבועי אין שורשים ממשיים. פרקו לשברים חלקיים מהצורה $\frac{A}{x}+\frac{Bx+C}{x^2+2x+5}$.
\end{hint}

\begin{hint}
באינטגרל של $\frac{Bx+C}{x^2+2x+5}$ פצלו למונה שהוא נגזרת המכנה ועוד קבוע, והשלימו לריבוע: $x^2+2x+5=(x+1)^2+4$.
\end{hint}

\begin{solution}
\textbf{פירוק המכנה.} $x^3+2x^2+5x=x(x^2+2x+5)$, והדיסקרימיננטה של $x^2+2x+5$ היא $4-20<0$, כך שהגורם הריבועי אי-פריק. מעלת המונה קטנה ממעלת המכנה, לכן ניתן לפרק ישירות לשברים חלקיים:
\[ \frac{4x^2+9x+10}{x(x^2+2x+5)}=\frac{A}{x}+\frac{Bx+C}{x^2+2x+5}. \]
כפל במכנה: $4x^2+9x+10=A(x^2+2x+5)+(Bx+C)x$.
\begin{itemize}
  \item הצבת $x=0$: $10=5A$, כלומר $A=2$.
  \item השוואת מקדמי $x^2$: $4=A+B$, כלומר $B=2$.
  \item השוואת מקדמי $x$: $9=2A+C$, כלומר $C=5$.
\end{itemize}
לכן
\[ f(x)=\frac{2}{x}+\frac{2x+5}{x^2+2x+5}. \]

\textbf{אינטגרציה.} $\int\frac{2}{x}\,dx=2\ln|x|$. בשבר השני נפצל $2x+5=(2x+2)+3$, כאשר $2x+2$ היא נגזרת המכנה:
\[ \int\frac{2x+5}{x^2+2x+5}\,dx=\int\frac{2x+2}{x^2+2x+5}\,dx+3\int\frac{dx}{(x+1)^2+4}. \]
האינטגרל הראשון שווה $\ln(x^2+2x+5)$ (הצבה $u=x^2+2x+5$; ללא ערך מוחלט כי הביטוי חיובי). בשני נציב $x+1=2t$:
\[ 3\int\frac{dx}{(x+1)^2+4}=\frac32\arctan\frac{x+1}{2}. \]

\textbf{תשובה.}
\[ \boxed{\int f(x)\,dx = 2\ln|x|+\ln(x^2+2x+5)+\frac32\arctan\frac{x+1}{2}+C} \]
(בכל אחד מהקטעים $x>0$ ו-$x<0$ בנפרד; ניתן לבדוק בגזירה שמתקבלת $f$.)
\end{solution}
